\chapter{Teoretická východiska}

\section{Anatomie ruky a horní končetiny}

Horní končetina (membrum superius) sestává z pletence a volné končetiny, která se dělí na tři typické úseky: stylopodium, zeugopodium a autopodium \cite{Cihak2011}. Autopodium je konečný úsek končetiny tvořený větším počtem skeletních elementů sestavených v paprsky, přičemž u horní končetiny je tento úsek přizpůsoben uchopovací funkci \cite{Cihak2011}. Kostru vlastní ruky tvoří celkem 27 kostí, které se dělí na kosti zápěstní, záprstní a články prstů \cite{Stuchla2024}. Osm zápěstních kostí (ossa carpi) je uspořádáno do proximální a distální řady a tvoří celek zvaný zápěstí \cite{Cihak2011}. Proximální řadu tvoří kosti loďkovitá, poloměsíčitá, trojhranná a hrášková, zatímco distální řada obsahuje kosti trapezium, trapezoideum, capitatum a hamatum \cite{Stuchla2024}. Uspořádání těchto kostí vytváří na dlaňové straně žlábek sulcus carpi, který je zesílen okrajovými vyvýšeninami \cite{Cihak2011}. Tento žlábek je uzavřen pevným vazem retinaculum musculorum flexorum, čímž vzniká karpální tunel (canalis carpi), kterým procházejí šlachy ohýbačů a nervus medianus \cite{Cihak2011}.
\paragraph{}Záprstí (metacarpus) tvoří pět kostí, z nichž první metakarp náležející k palci je nejkratší a nejsilnější \cite{Cihak2011}. Prsty mají celkem 14 článků, přičemž palec disponuje dvěma články a ostatní prsty mají články tři \cite{Stuchla2024}. Klíčovým mechanismem pro úchop je opozice palce, kterou umožňuje sedlový kloub mezi os trapezium a prvním metakarpem \cite{Cihak2011}. Krátké svaly ruky jsou rozděleny do skupin thenaru, hypothenaru, lumbrikálních a mezikostních svalů \cite{Stuchla2024}. Skupina thenaru na palcové straně zahrnuje krátký odtahovač, krátký ohýbač, oponující sval a přitahovač palce \cite{Cihak2011}. Čtyři červovité svaly (mm. lumbricales) začínají v dlani na šlachách hlubokého ohýbače prstů \cite{Cihak2011}. Mezikostní svaly jsou uloženy v intermetakarpálních prostorech a dělí se na tři palmární a čtyři dorsální svaly rozložené podle osy jdoucí třetím prstem \cite{Cihak2011}. Dlaň je vyztužena tuhou vazivovou ploténkou aponeurosis palmaris, která je pevně spojena s kůží \cite{Cihak2011}. Periferní nervy ruky jsou tvořeny svazky vláken obklopených endoneuriem, perineuriem a epineuriem \cite{Cihak2016}.

\section{Fyziologie úchopu}

Vývoj horní končetiny začíná kolem 24. dne embryonálního života vznikem končetinového pupenu v úrovni somitů C5 až Th1. Nositlem morfogenetické informace o tvaru a poloze svalu je vazivový mesenchym, nikoliv samotné svalové buňky. Během růstu se horní končetina v oblasti paže otáčí zevně a předloktí do částečné pronace, což umožňuje výsledné funkční postavení ruky \cite{Cihak2011}. Mezi základní funkční vlastnosti svaloviny ruky patří dráždivost, kontraktibilita, protažitelnost a elasticita. Signálem pro svalovou kontrakci je uvolnění acetylcholinu na nervosvalové ploténce, což aktivuje sarkoplazmatické retikulum k vyplavení vápenatých iontů \cite{Stuchla2024}. Energie pro tento proces je získávána štěpením molekul adenosintrifosfátu (ATP) \cite{Cihak2011}, přičemž svaly ruky fungují jako motorické orgány měnící energii chemických vazeb v teplo a mechanickou práci. Podstatou stahu při úchopu je dle Huxleyovy a Hansenovy teorie relativní pohyb (zasouvání) vláken aktinu mezi vlákna myozinu \cite{Benes2022}.
Po nadprahovém podráždění svalové napětí rychle roste k maximu během 50 ms a následně během 100–150 ms klesá. Celková uvolněná energie se skládá z aktivačního tepla, zkracovacího tepla a mechanické práce, přičemž účinnost této přeměny dosahuje maximálně 40 \%. Izometrická kontrakce mění pouze napětí svalu bez změny délky, sval při ní nevykonává mechanickou práci, slouží k fixaci předmětu a vyvíjí při ní největší sílu \cite{Benes2022}, zatímco izotonická kontrakce mění délku svalu při vykonávání pohybu, slouží ke zrychlení a vykonává kladnou mechanickou práci \cite{Benes2022, Stuchla2024}. Pokud je vnější síla působící na ruku velká, dochází k excentrické činnosti (prodloužení), kdy sval pohyb zpomaluje a vykonává zápornou práci \cite{Benes2022}.
\paragraph{}Lidský sval je schopen vyvinout silové působení přibližně 40 N·cm⁻², které při úchopu v těle působí nejčastěji prostřednictvím jednozvratných pák. Rychlost zkracování svalu je největší na počátku pohybu a postupně klesá k nule v okamžiku maximálního zkrácení \cite{Benes2022}. Hmatové vnímání zajišťují mechanoreceptory, které přeměňují mechanickou energii na akční potenciály \cite{Cihak2016}. Meissnerova tělíska v bříškách prstů jsou rychle se adaptující receptory citlivé na jemné nízkofrekvenční chvění. Merkelovy disky informují o trvajícím kontaktu epidermis s předmětem a umožňují tak vnímání tlaku. Vaterova-Paciniho tělíska v podkoží reagují na vysokofrekvenční vibrace a představují nejsložitější kožní mechanoreceptory. Termoreceptory reagující na teplo vykazují maximální aktivitu při 45°C, zatímco chladové receptory jsou nejaktivnější při 28°C. Nociceptory vnímající bolest jsou na rozdíl od hmatových receptorů neadaptivní a reagují na podněty poškozující tkáň \cite{Stuchla2024}.
\paragraph{}Koordinace jemných a přesných pohybů prstů je řízena korovým mozečkem. Hlubo -    ké čití (propriocepci) zajišťují svalová vřeténka informující o délce svalu a Golgiho šlachová tělíska detekující napětí šlach \cite{Stuchla2024}. Nervus medianus je pro úchopovou funkci zásadní, neboť jeho poškození znemožňuje opozici palce a ohýbání druhého a třetího prstu. Pokud dojde k přerušení nervu, probíhá antegrádně od místa poškození proces rozpadu axonu a myelinové pochvy zvaný Wallerova degenerace. Regenerace nervových vláken v ruce je možná pouze za přítomnosti Schwannových buněk, které slouží jako vodivé struktury. Podnět je v nervovém vláknu vždy překódován do specifické frekvence akčních potenciálů \cite{Cihak2016}.


\section{Biomechanika úchopu a silové interakce}

Biomechanika ruky a úchopu představuje komplexní studium vnitřních a vnějších sil působících na pohybový aparát horní končetiny. Základními stavebními prvky tkání jsou vlákna elastinu, která umožňují pružné deformace až o 150 \%, a tuhá kolagenní vlákna zajišťující pevnost v tahu. Mechanické chování kostí se při malém zatížení řídí Hookovým zákonem, podle něhož je deformace přímo úměrná působící síle. Pro natažení nebo stlačení kosti o délce $l$ a průřezu $S$ platí vztah:
\begin{equation}
\Delta l = \frac{l \cdot F}{E \cdot S}
\end{equation}
kde $E$ představuje Youngův modul pružnosti. Pevnost dlouhých kostí v ohybu je zvýšena jejich trubicovitou strukturou, přičemž výpočet průhybu $s$ využívá rozdílu čtvrtých mocnin vnějšího a vnitřního poloměru trubice ve vzorci:
\begin{equation}
s = \frac{F \cdot l^3}{12 \cdot \pi \cdot E \cdot (R_2^4 - R_1^4)}
\end{equation} 
Struktura kostí se navíc dynamicky přizpůsobuje zátěži podle Wolfova zákona o transformaci kostí \cite{Benes2022}.

\paragraph{}Vlastní mechanickou práci úchopu vykonávají svaly, v nichž se chemická energie mění v teplo a pohyb na základě Huxleyovy a Hansenovy teorie o zasouvání aktinových a myozinových filament \cite{Benes2022}. Maximální izometrická aktivní síla svalu $F_{0M}$ je definována jako součin specifického napětí $\sigma$ a fyziologického průřezu svalu $PCSA$. Fyziologický průřez lze vypočítat ze vztahu:
\begin{equation}
PCSA = \frac{m \cdot \cos(\alpha)}{\rho \cdot L_{0M}}
\end{equation}
kde $m$ je hmotnost svalu, $\alpha$ úhel zpeření, $\rho$ hustota a $L_{0M}$ optimální délka svalu. Celková energie $E$ uvolněná při svalovém stahu se dělí na aktivační teplo $W_a$, zkracovací teplo $W_z$ a vykonanou mechanickou práci $W$ podle rovnice $E = W_a + W_z + W$. Účinnost této přeměny v lidském těle dosahuje maximálně 40 \% \cite{Capek2018}.

\paragraph{}Vztah mezi silou a rychlostí svalové kontrakce popsal A. V. Hill hyperbolickou závislostí ve tvaru:
\begin{equation}
F_{0M} \cdot f_v(v_M) = F_{0M} \cdot \frac{v_{0M} - v_M}{v_{0M} + c \cdot v_M}
\end{equation}
kde $v_M$ je rychlost zkracování a $v_{0M}$ maximální rychlost kontrakce. Aktivní složka svalové síly v závislosti na délce svalu vykazuje parabolický průběh popsaný vztahem:
\begin{equation}
F_{La} = F_{0M} \cdot \left(1 - \left(\frac{L_M}{L_{0M}} - 1\right) \cdot 2\right)^2
\end{equation}
Pasivní odpor svalu při protažení naopak roste exponenciálně podle funkce:
\begin{equation}
F_{pe} = F_{0M} \cdot \left(\frac{L_M}{L_{0M}}\right)^3 \cdot \exp\left(8 \cdot \frac{L_M}{L_{0M}} - 12.9\right)
\end{equation}
Celková síla svalu $F_M$ v dynamickém režimu je pak součtem pasivních sil, aktivních sil a viskózního členu tlumícího pohyb. Síla je na kostru přenášena přes šlachy, které se chovají jako exponenciální pružiny s deformací $\varepsilon_T = (L_T - L_{sT}) / L_{sT}$, kde $L_{sT}$ je volná délka šlachy \cite{Capek2018}.

\paragraph{}Řízení úchopu vychází z elektrické aktivity nervové soustavy, kde základní membránové napětí $U$ určuje Nernstova rovnice \cite{Mornstein2018}:
\begin{equation}
U = \frac{RT}{zF} \cdot \ln\left(\frac{c}{c^*}\right)
\end{equation}
Pro dosažení rovnováhy na obou stranách membrány s ohledem na různé ionty se využívá Donnanova rovnováha \cite{Mornstein2018} ve tvaru:
\begin{equation}
-U = \frac{RT}{F} \cdot \ln\left(\frac{c_{K^+}}{c_{K^+}^*}\right) = \frac{RT}{F} \cdot \ln\left(\frac{c_{Cl^-}}{c_{Cl^-}^*}\right)
\end{equation}
Proces svalové aktivace $a(t)$ v závislosti na nervovém dráždění $u(t)$ lze modelovat diferenciální rovnicí prvního řádu \cite{Capek2018}:
\begin{equation}
\dot{a} = \frac{u^2 - u \cdot a}{\tau_{rise}} + \frac{u - a}{\tau_{fall}}
\end{equation}
Lidská ruka disponuje 22 stupni volnosti \cite{Brazdil2020}, přičemž zápěstí přidává další dvě. U celkového biomechanického modelu lidské ruky bychom tedy mluvili o 24 DOF \cite{Hansen2018}. Pohyblivost mechanismu prstu $i$ lze ověřit kinematickým rozborem $i = 3(n - 1) - \sum \xi_i$ \cite{Safran2022}.

\paragraph{}Při využití technických systémů pro podporu úchopu, jako jsou myoelektrické protézy, se k vyhodnocení otřesů používá výpočet euklidovské vzdálenosti zrychlení:
\begin{equation}
j_{area}(j) = \sqrt{X_i^2 + Y_i^2 + Z_i^2}
\end{equation}
Úspěšnost automatické klasifikace hmatových gest se měří pomocí přesnosti:
\begin{equation}
ACC = \frac{TP + TN}{TP + FP + TN + FN}
\end{equation}
Pro jemné rozlišení mezi jednotlivými typy úchopů se v algoritmech SVM často využívá RBF kernel definovaný vztahem \cite{Brazdil2020}:
\begin{equation}
K(x_j, x_k) = \exp(-\gamma \cdot \|x_j - x_k\|^2)
\end{equation}
Přesnost silové interakce je zajištěna zpětnou vazbou z mechanoreceptorů, jako jsou Meissnerova tělíska reagující na chvění nebo Merkelovy disky vnímající trvalý tlak. Samotná síla úchopu je v nervových vláknech kódována frekvencí akčních potenciálů, která je úměrná velikosti mechanického podnětu \cite{Stuchla2024}.

\subsection{Typy úchopů a přenos síly}

V robotických a bionických systémech je problematika úchopů a přenosu síly ústředním tématem, které propojuje biologické principy s pokročilým inženýrstvím. Základním stavebním kamenem je lidský somatosenzorický systém, který využívá síť mechareceptorů (jako jsou Ruffiniho tělíska, Merkelovy disky či Paciniho tělíska) k detekci mechanických změn, vibrací a protažení kůže \cite{Cihak2016}. Tato biologická inspirace se promítá do technického řešení tzv. paradigmatu Leader-Follower (Vůdce-Následovník), kde operátor skrze haptické rozhraní (Leader) ovládá vzdálený robotický systém (Follower), který interaguje s prostředím \cite{GIUDICI2025}. Cílem těchto systémů je dosáhnout vysoké míry transparence, kdy operátor vnímá interakci s objektem tak, jako by se ho dotýkal přímo \cite{GIUDICI2025}.
\subsubsection{Typologie a klasifikace úchopů}
V oblasti bioniky a robotické telemanipulace se úchop definuje jako statická poloha ruky umožňující bezpečné držení objektu \cite{Nyrobtseva2018}. Úchopy lze klasifikovat podle geometrie a počtu kontaktních bodů na bodové (např. tříprsté) a osové (všemi prsty), přičemž osové úchopy vykazují vyšší stabilitu díky schopnosti lépe detekovat rotační a spinové tlaky \cite{Trajlinek2021}. Mezi základní implementované typy patří špetkový úchop (pinch grip), realizovaný přitlačením špičky palce k ukazováčku, a válcový úchop (cylindrical grip) určený pro manipulaci s rotačními tělesy \cite{Nyrobtseva2018}. Výzkumy také definují kategorii „hákového“ úchopu (hook) pro přenášení břemen \cite{Nyrobtseva2018, GIUDICI2025}.
Pro účely kontinuálního rozpoznávání pohybů se zavádí koncept zakřivení prstů (finger curvature), který zjednodušuje komplexní anatomii ruky na proměnnou v rozsahu \cite{GIUDICI2025}, což snižuje počet stupňů volnosti pro zpracování řídicími algoritmy \cite{Sun2023}. Manipulační úlohy se dále dělí na kategorie jako re-grasping (přehmátnutí), finger gaiting (střídavý pohyb prstů) nebo manipulaci uvnitř úchopu (in-grasp manipulation) \cite{GIUDICI2025}. Distribuce síly při těchto úkonech je specifická; například u dvouprstého uchopování se hlavní tlak soustřeďuje do oblasti palce a ukazováčku \cite{Nyrobtseva2018}.
\subsubsection{Mechanismus přenosu a snímání síly}
Efektivní přenos síly vyžaduje přesnou detekci kontaktu na straně Followera. Nejrozšířenější technologií jsou silově citlivé rezistory (FSR), které mění svůj elektrický odpor v závislosti na mechanickém tlaku \cite{Nyrobtseva2018, Trajlinek2021}. Tyto senzory bývají integrovány do bříšek prstů (fingertips) a umožňují monitorovat intenzitu stisku v reálném čase \cite{Trajlinek2021}. Pokročilejší systémy využívají magnetické taktilní senzory založené na Hallově jevu. Ty měří deformaci silikonového těla senzoru skrze změny magnetického pole, což umožňuje detekovat nejen normálové, ale i smykové síly \cite{GIUDICI2025}. Pro odhad svalové síly uživatele se využívá povrchová elektromyografie (sEMG), kde signály z předloktí interpretují neurální sítě typu LSTM pro přesný odhad motorického záměru \cite{Sun2023}.
\subsubsection{Haptická zpětná vazba a energetika úchopu}
Naměřená síla je přenášena zpět k uživateli skrze různé modality zpětné vazby. Kinestetická zpětná vazba využívá exoskeletální rukavice (např. HGlove) k vyvolání reálného odporu v kloubech prstů operátora, což umožňuje vnímání tuhosti (stiffness) vzdáleného objektu \cite{GIUDICI2025}. Alternativou je vibrotaktilní vazba, kde miniaturní vibrační motorky v rukavici signalizují moment prvního kontaktu nebo intenzitu tlaku \cite{Papakonstantinou20251218, Trajlinek2021}.
Implementace taktilní vazby má přímý dopad na energetickou účinnost robotického systému. V přítomnosti haptické odezvy operátor dokáže identifikovat optimální bod úchopu, čímž se vyhne nadbytečnému krouticímu momentu motorů a výrazně sníží spotřebu energie \cite{Papakonstantinou20251218}. Bezpečnost operátora je zajištěna programovými limity, kdy se zpětnovazebná síla často omezuje na bezpečnou hodnotu (např. 5 N). Pro pokročilé aplikace v prostředí bez vizuální kontroly lze taktilní data využít i k realitní 3D rekonstrukci tvaru objektu ve virtuální realitě \cite{GIUDICI2025}.

\subsection{Kontakt prst–objekt a percepce síly}
Tady je sloučený text:

Fyzikální podstatu kontaktu mezi prstem a objektem definuje vzájemné působení normálové síly, působící kolmo na povrch, a tečné (smykové) síly vyplývající z tření \cite{Dutta2002, GIUDICI2025}. Stabilita úchopu vyžaduje, aby sevřívací síla přesáhla minimální hodnotu potřebnou k udržení objektu, násobenou příslušným bezpečnostním faktorem \cite{Mayer2020427}, přičemž závisí rovněž na tvaru a poloze předmětu a na geometrii uchopovače \cite{Devaraja20201222}. Kritickým jevem je skluz (\textit{slippage}), k němuž dochází v okamžiku, kdy smykové síly překonají statické tření; tomuto jevu obvykle předchází série mikroskluzů detekovatelných na styčné ploše \cite{Begalinova2020, Xu2025}. Detekce počínajících mikroskluzů je proto klíčová pro včasnou korekci úchopové síly a prevenci pádu předmětu.
Biologické vnímání kontaktní síly v bříšcích prstů zajišťuje somatosenzorický systém, jehož základ tvoří čtyři typy specializovaných mechanoreceptorů lišících se frekvenční citlivostí a rychlostí adaptace \cite{Stuchla2024, Cihak2016}. Meissnerova tělíska (FA~I) jsou rychle se adaptující receptory citlivé na nízkofrekvenční vibrace v rozsahu 30–50~Hz; detekují texturu a počátek skluzu a tvoří přibližně 40~\% receptorů na ruce \cite{Stuchla2024, Cihak2016}. Merkelovy disky (SA~I) jsou pomalu se adaptující receptory reagující na trvalý tlak a lehký dotek, klíčové pro vnímání hran a hrubých textur, s podílem přibližně 25~\% \cite{Stuchla2024, Cihak2016}. Paciniho tělíska (FA~II) jsou velmi rychle se adaptující receptory uložené hluboko v tkáni, citlivé na vysokofrekvenční vibrace v rozsahu 250–300~Hz a náhlé změny \cite{Stuchla2024, Cihak2016}. Ruffiniho zakončení (SA~II), tvořící přibližně 20~\% receptorů, se pomalu adaptují a detekují napínání kůže při ohybu kloubů a silovém úchopu \cite{Stuchla2024, Cihak2016}. Prahové hodnoty lidského hmatu jsou pozoruhodně nízké: na dlani je člověk schopen zaznamenat hmotnost od 0,019~g u žen do 0,055~g u mužů, přičemž práh pro vnímání jemného doteku se pohybuje v rozmezí 6 až 20~g \cite{Stuchla2024}.
Z hlediska technického snímání kontaktní síly se v literatuře nejčastěji uvádějí tři přístupy. Odporové senzory citlivé na sílu (FSR) vynikají nízkou cenou, malou tloušťkou, odolností vůči elektromagnetickému rušení a snadnou integrací do konečků prstů, avšak jejich nevýhodou je nelineární odezva, hystereze a nižší přesnost při měření absolutní velikosti síly \cite{Sadun2016711, Lebosse2011}. Alternativu představuje nepřímá detekce prostřednictvím monitorování proudu servomotorů, která nevyžaduje žádné senzory v místě kontaktu a je vhodná pro odhad celkové úchopové síly; mechanické tření převodů a šum v odběru proudu však maskují jemný kontakt a znemožňují citlivou detekci prvotního doteku \cite{Mohammadi2024418, Ninu2014}. Moderní taktilní senzory -- například pole na bázi Hallova jevu -- umožňují měřit sílu ve třech osách a poskytují detailní prostorovou informaci o rozložení tlaku včetně smykových složek, avšak za cenu vyšší pořizovací ceny, hardwarové objemnosti a výpočetní náročnosti zpracování dat \cite{GIUDICI2025, Begalinova2020, Xu2025}.
Skluz je v literatuře chápán jako makroskopická událost předcházená sérií mikroskluzů \cite{Begalinova2020}. Jeho detekce je nezbytná pro uzavření regulační smyčky: jakmile systém zaznamená pohyb objektu v úchopu, musí automaticky zvýšit úchopovou sílu, aby zabránil jeho vypadnutí \cite{Begalinova2020, Xu2025}. K detekci se využívá extrakce vibračních příznaků, analýza změn kontaktních sil nebo pokročilé algoritmy strojového učení, jako jsou architektury CNN-LSTM umožňující predikci v reálném čase \cite{Xu2025, GIUDICI2025}.
Vnímání kontaktní síly úzce souvisí se stabilitou robotického úchopu jako celku. Přenos informace o síle k operátorovi umožňuje vnímat okamžik kontaktu, což vede k jemnější regulaci tlaku, snížení energetické náročnosti a rychlejšímu dosažení stabilního úchopu \cite{Papakonstantinou20251218, Thomas2023110}. Bez silové zpětné vazby je manipulace s křehkými či deformovatelnými předměty obtížná, neboť hrozí jejich poškození nebo vyklouznutí \cite{Park2022, DAbbraccio201923, GIUDICI2025}. Pro zajištění stability při práci s neznámými objekty se doporučuje senzorová fúze kombinující FSR senzory pro rychlou detekci kontaktu s měřením proudu motorů pro přesnější určení tuhosti a stability v čase \cite{Mohammadi2024418}. Důležitým faktorem je rovněž latence systému: překročí-li zpoždění hodnotu 300–400~ms, uživatel vnímá pohyb jako neplynulý, což narušuje schopnost udržet stabilní interakci s objektem \cite{Bolarinwa2022}.

\section{Principy haptiky a silové zpětné vazby}
Haptika představuje vědní obor zabývající se hmatovou interakcí a simulací fyzického kontaktu, což je v kontextu moderní teleoperace a robotických systémů klíčové pro dosažení vysoké míry imerze a přesnosti dálkového ovládání \cite{GIUDICI2025, Patel2022}. Zatímco somatosenzorický systém člověka přirozeně integruje komplexní vjemy o tlaku, vibracích a pohybu, v bilaterálních systémech konfigurace Leader–Follower musí být tato informace uměle zprostředkována, aby operátor mohl intuitivně vnímat fyzikální vlastnosti vzdálených objektů a adekvátně na ně reagovat včasnou úpravou úchopové síly či trajektorie pohybu \cite{Patel2022, Thomas2023110, GIUDICI2025}. Je přitom nezbytné odlišovat haptiku jako komplexní interdisciplinární disciplínu od její konkrétní technické realizace, která prostřednictvím mechanických, elektrických či vibračních aktuátorů aktivně stimuluje lidské receptory a převádí digitální data zpět na fyzikální podnět \cite{Yang2025, Li2025, Patel2022}.
V rámci širokého spektra haptických modalit zaujímá specifické postavení silová zpětná vazba, označovaná jako \textit{force feedback}, která cílí na kinestetické vnímání interakčních sil prostřednictvím receptorů ve svalech, šlachách a kloubech aktivním odporem kladeným proti pohybu \cite{Patel2022, Yang2025}. Na rozdíl od taktilních (kožních) vjemů informujících o textuře či vibracích umožňuje silová zpětná vazba uživateli fyzicky pocítit tuhost, hmotnost a prostorový odpor manipulovaného objektu, čímž zásadně zvyšuje stabilitu a bezpečnost interakce zejména v náročných teleoperačních scénářích \cite{GIUDICI2025, Patel2022, Li2025}. Návrh těchto systémů vyžaduje precizní zvládnutí regulačních smyček a potlačení nežádoucích jevů; následující podkapitoly se proto detailně zaměří na systematickou klasifikaci haptických modalit, specifická omezení silové zpětné vazby plynoucí z mechanických ztrát a komunikačních latencí a na kritickou otázku stability haptických systémů v uzavřeném teleoperačním řetězci.
\subsection{Klasifikace haptické zpětné vazby}
Haptická percepce a její technická simulace se v~odborné literatuře systematicky dělí na dvě hlavní složky: kožní (taktilní) a kinestetickou (proprioceptivní) \cite{Yang2025, GIUDICI2025, Li2025}.
Taktilní vnímání je zprostředkováno receptory v~pokožce, které detekují jemný tlak, texturu a vibrace, zatímco kinestetická zpětná vazba cílí na receptory v~kloubech, svalech a šlachách, čímž uživateli umožňuje vnímat polohu končetin, hmotnost a tuhost manipulovaných předmětů \cite{Yang2025, Belcamino2024, DAbbraccio201923}.

Mezi nejrozšířenější metody taktilní stimulace patří vibrotaktilní zpětná vazba využívající vibrační motorky s~excentrickou hmotou (ERM) nebo lineární rezonanční aktuátory (LRA), které jsou vysoce nositelné a energeticky úsporné, avšak neposkytují vjem fyzického odporu a trpí senzorickou adaptací při déletrvajícím působení \cite{Kourtesis202271, Park2022, Thomas2023110}.
Alternativou je elektrotaktilní stimulace, která pomocí elektrod dráždí nervová zakončení přímo na povrchu kůže, čímž umožňuje konstrukci velmi kompaktních zařízení bez pohyblivých částí \cite{Kourtesis202271, Yang2025}.
Nevýhodou je nelineární charakter vjemu a nutnost individuální kalibrace z~důvodu proměnné impedance kůže.
Chování membrány nervového vlákna při elektrotaktilní stimulaci lze popsat diferenciální rovnicí ve tvaru
%
\begin{equation}
    -G_a \frac{\mathrm{d}^2}{\mathrm{d}x^2} V_m
    + C_m \frac{\mathrm{d}}{\mathrm{d}t} V_m
    + C_m V_m
    = G_a \frac{\mathrm{d}^2}{\mathrm{d}x^2} \Psi,
\end{equation}
%
kde $G_a$~je vnitřní vodivost axonu, $V_m$~membránové napětí a $\Psi$~vnější potenciál stimulačního signálu \cite{Yang2025}.

Z~hlediska věrnosti vjemu a vhodnosti pro náročné teleoperační úlohy vyžadující přesnou regulaci úchopu je za nejúčinnější považována kinestetická silová zpětná vazba realizovaná nositelnými exoskeletálními systémy, jako je rukavice \textit{Remote Feelings} \cite{Li2025, GIUDICI2025, Giudici2024}.
Tato zařízení využívají aktivní mechanické aktuátory -- nejčastěji DC~motory nebo lanková vedení -- k~vyvolání odporu přímo proti pohybu prstů, čímž věrně simulují mechanickou impedanci vzdáleného prostředí \cite{Li2025, GIUDICI2025}.
Přestože jsou tyto mechanické systémy konstrukčně nejsložitější a typicky méně nositelné než vibrační či elektrické modality, poskytují operátorovi nezbytnou informaci o~okamžiku kontaktu a tuhosti prostředí v~reálném čase, což zásadně zvyšuje stabilitu dálkového ovládání a úspěšnost manipulace s~objekty \cite{Li2025, Park2022, Giudici2024}.
\subsection{Force feedback a jeho omezení}
Silová zpětná vazba, odborně definovaná jako kinestetická haptika, zprostředkovává operátorovi teleoperačního systému informace o~interakčních silách prostřednictvím receptorů ve svalech, šlachách a kloubech \cite{Patel2022, Li2025}.
Tato modalita se zásadně liší od taktilní (kožní) vazby, která stimuluje receptory v~pokožce pro vnímání textury, vibrací či lokálního tlaku \cite{Patel2022, Yang2025}.
V~kontextu dálkového ovládání robotů je silová odezva klíčová pro navrácení chybějícího smyslového vjemu o~tuhosti a hmotnosti vzdáleného prostředí, což umožňuje intuitivní regulaci úchopové síly, prevenci poškození křehkých objektů a bezpečnou manipulaci \cite{GIUDICI2025, Patel2022, Park2022}.

Ústředním cílem návrhu těchto systémů je transparence, tedy stav, kdy operátor vnímá vzdálenou úlohu, jako by s~ní interagoval přímo.
Bilaterální teleoperaci lze matematicky vyjádřit pomocí hybridní matice $\mathbf{H}(s)$ vztahující síly ($f$) a rychlosti ($v$) na straně Leader a Follower:
%
\begin{equation}
    \begin{bmatrix} f_h \\ -v_n \end{bmatrix} =
    \begin{bmatrix} h_{11} & h_{12} \\ h_{21} & h_{22} \end{bmatrix}
    \begin{bmatrix} v_o \\ f_n \end{bmatrix},
\end{equation}
%
kde v~ideálně transparentním systému nabývají prvky matice hodnot $h_{11} = h_{22} = 0$, $h_{12} = 1$ a $h_{21} = -1$, což značí dokonalou shodu sil a pohybů bez vlivu parazitní impedance samotného technického zařízení \cite{Black2025, Almeida202098}.
V~praxi je však transparence omezena mechanickými ztrátami, jako je tření a hystereze v~převodech, šířkou pásma aktuátorů a především latencí komunikačního řetězce \cite{Patel2022, Black2025}.

Zpoždění přenosu dat vnáší do uzavřené smyčky fázový posun, který může způsobit nestabilitu projevující se oscilacemi; již latence v~řádu 50~ms může u~operátora významně zkreslit percepci tuhosti objektu \cite{Patel2022}.
Zatímco zpoždění do přibližně 300~ms umožňuje kontinuální řízení, latence překračující 400~ms nutí uživatele přecházet na diskontinuální strategii ovládání s~výraznými dopady na plynulost a intuitivnost pohybu \cite{Bolarinwa2022, Almeida202098}.
Bezdrátový přenos tyto problémy dále prohlubuje vlivem proměnného zpoždění (jitteru) a výpadků paketů; přičemž průměrná hodnota latence je pro kvalitu řízení obvykle kritičtější než samotná nahodilost jejího kolísání \cite{Suman2025, Almeida202098}.
Konkrétní srovnání ukázalo, že protokol Bluetooth~Classic vykazuje vyšší latenci a nižší propustnost než alternativní protokoly, jako je ESP-NOW, což vede k~horší synchronizaci haptického vjemu s~vizuální odezvou v~reálném čase \cite{Eridani2021918, Papakonstantinou20251218}.

Kvalitu zpětné vazby dále limituje nepřímé odvozování síly z~proudového odběru motorů.
Tato metoda je konstrukčně jednodušší a nevyžaduje senzory přímo v~místě kontaktu, avšak vnitřní odpor a mechanické tření motoru generují šum maskující jemné změny při prvotním dotyku, což omezuje citlivost a přesnost vnímání odporu prostředí \cite{Patel2022, Mohammadi2024418}.
Ke kompenzaci tohoto omezení literatura doporučuje senzorovou fúzi proudového monitoringu se senzory přímého tlaku, která zvyšuje citlivost detekce v~počáteční fázi kontaktu a umožňuje nastavit bezpečné silové limity zabraňující poškození jak manipulovaných objektů, tak samotného zařízení \cite{Lebosse2011, Mohammadi2024418}.
\subsection{Stabilita haptických systémů}
Stabilita bilaterálních teleoperačních systémů se silovou zpětnou vazbou je kriticky ovlivněna akumulací energie v~uzavřené regulační smyčce, což je jev způsobený především zpožděním přenosu dat mezi stranami Leader a~Follower \cite{Patel2022, Black2025}.
V~těchto systémech vnáší latence do zpětné vazby fázový posun vedoucí k~nekontrolovaným oscilacím, jež operátor vnímá jako trhavý pohyb nebo efekt \textit{chatteringu} \cite{Patel2022, Black2025, GIUDICI2025}.
Klíčovými parametry ovlivňujícími stabilitu jsou tuhost vzdáleného prostředí ($Z_e$), dynamika operátora ($Z_o$) a~komunikační latence ($\tau$) \cite{Black2025}.
Již latence kolem 50~ms může významně zkreslit percepci tuhosti objektu; hodnoty do přibližně 300~ms zpravidla umožňují kontinuální řízení, zatímco zpoždění nad 500~ms nutí operátora přejít na diskontinuální strategii ovládání \cite{Patel2022, Almeida202098, Bolarinwa2022}.

Fundamentálním nástrojem pro zajištění stability nezávislé na zpoždění je princip pasivity (\textit{passivity-based control}), který koncipuje teleoperační systém jako pasivní dvouportovou síť negenerující vlastní energii \cite{Patel2022, Black2025}.
Pasivita systému s~hybridní maticí $\mathbf{H}(s)$ je formálně ekvivalentní podmínce, že jeho rozptylová (\textit{scattering}) matice
%
\begin{equation}
    \mathbf{S}(s)
    = \bigl(\mathbf{H}(s) - b\,\mathbf{I}\bigr)
      \bigl(\mathbf{H}(s) + b\,\mathbf{I}\bigr)^{-1},
\end{equation}
%
kde $b > 0$ je vlnová impedance a $\mathbf{I}$ jednotková matice, splňuje $\|\mathbf{S}(j\omega)\|_2 \leq 1$ pro všechna $\omega$ \cite{Black2025, Patel2022}.
Praktická implementace tohoto principu nejčastěji využívá metodu Wave Variable Control (WVC), která transformuje přenášené proměnné tak, aby přenosový kanál zůstal pasivní i~při proměnném zpoždění, nebo Time-Domain Passivity Approach (TDPA), která monitoruje tok energie v~reálném čase a~v~případě detekce nepasivního chování vnáší do systému adaptivní tlumení \cite{Patel2022, Bolarinwa2022}.

Bezdrátový přenos vnáší do stability specifická rizika v~podobě kolísavé latence (jitteru) a~výpadků paketů, přičemž průměrná hodnota zpoždění je pro integritu zpětné vazby kritičtější než samotná nahodilost jejího kolísání \cite{Suman2025, Aijaz2018}.
Srovnání protokolů ukázalo, že Bluetooth~Classic vykazuje vyšší latenci a~nižší spolehlivost datového toku než alternativní protokoly (např.\ ESP-NOW), což se projevuje nespojitými skoky v~silové odezvě a~následnou destabilizací uzavřené smyčky \cite{Eridani2021918, Suman2025}.

Pro praktickou implementaci literatura doporučuje víceúrovňové bezpečnostní mechanismy \cite{Patel2022, Black2025, Bolarinwa2022}.
Shora omezené hodnoty akčních zásahů motorů (saturace PWM) chrání před mechanickým poškozením zařízení i~manipulovaného objektu.
Softwarová hystereze s~průměrováním vzorků kontaktního signálu eliminuje vliv šumu a~zabraňuje falešným aktivacím zpětné vazby.
Watchdog mechanismy a~timeouty při detekci ztráty spojení nebo překročení kritického zpoždění automaticky deaktivují silovou zpětnou vazbu a~uvádějí systém do definovaného bezpečného stavu.
Účinným doplňkem je rovněž nízkofrekvenční filtrace zpětné vazby, která injektuje do řetězce virtuální tlumení a~potlačuje oscilace vznikající při kontaktu s~tuhými objekty \cite{Black2025, Patel2022}.

\section{Povrchová elektromyografie (sEMG)}
Povrchová elektromyografie (sEMG) představuje neinvazivní metodu založenou na snímání a~analýze bioelektrických potenciálů, které vznikají při kontrakci kosterního svalstva a~jsou detekovány pomocí elektrod umístěných přímo na povrchu pokožky \cite{ALCAN20231026, Reaz2006, Wu2024613}.
V~kontextu moderních robotických a~protetických systémů slouží sEMG jako fundamentální řídicí modalita, která zprostředkovává přirozené rozhraní mezi fyziologickým záměrem uživatele a~následným mechanickým pohybem technického zařízení \cite{ALCAN20231026, Song2023823, Xiong2024}.
Fyziologickou podstatu tohoto signálu tvoří časoprostorová superpozice akčních potenciálů motorických jednotek (MUAP), kterou lze matematicky popsat modelem
%
\begin{equation}
    x(n) = \sum_{r=0}^{N-1} h(r)\,e(n-r) + w(n),
\end{equation}
%
kde $x(n)$ představuje modelovaný sEMG signál, $e(n)$ reprezentuje impulzy nervového buzení, $h(r)$ odpovídá odezvě motorické jednotky a~$w(n)$ zahrnuje aditivní bílý gaussovský šum \cite{Reaz2006}.
Zásadní výhodou sEMG oproti jiným modalitám je skutečnost, že signál se objevuje přibližně 50 až 100~ms před samotnou fyzickou aktivitou svalu, což umožňuje predikovat záměr pohybu v~reálném čase a~potenciálně kompenzovat latence akčních členů \cite{Song2023823, Xiong2024}.
Navzdory vysoké míře intuitivnosti však praktické nasazení naráží na výzvy v~podobě stochastického a~nestacionárního charakteru signálu, vysoké náchylnosti k~elektromagnetickému rušení, parazitních pohybových artefaktů a~značné variability mezi jednotlivými uživateli i~polohami elektrod \cite{Chowdhury2013917, DeLuca2010, Sul2025627, Song2023823}.
Následující podkapitoly se proto detailně zaměří na fyziologické mechanismy vzniku sEMG v~rámci motorických jednotek a~následně na specifické charakteristiky a~technická omezení, která definují stabilitu a~přesnost myoelektrického ovládání v~teleoperačním řetězci.
\subsection{Fyziologický původ sEMG signálu}
Povrchový elektromyografický signál má svůj fyziologický původ v~motorické jednotce, která je definována jako alfa-motoneuron a~všechna svalová vlákna jím inervovaná \cite{Reaz2006, Stuchla2024}.
Proces vzniku signálu je iniciován nervovým vzruchem, který po dosažení nervosvalové ploténky uvolňuje mediátor acetylcholin; ten vyvolá depolarizaci membrány svalového vlákna, a~tím generuje akční potenciál vedoucí k~mechanické kontrakci \cite{ALCAN20231026, Reaz2006, Stuchla2024}.
Výsledný sEMG signál snímaný na povrchu kůže nepředstavuje aktivitu jediného vlákna, nýbrž je výsledkem časoprostorové superpozice akčních potenciálů mnoha motorických jednotek (MUAP), což vytváří takzvaný interferenční obrazec \cite{ALCAN20231026, Chowdhury2013917, Reaz2006, Song2023823, Sul2025627}.

Úroveň svalové aktivace je v~signálu kódována prostřednictvím náboru nových motorických jednotek a~frekvence jejich impulsů (\textit{firing rate}); vztah mezi silou kontrakce a~amplitudou sEMG však není striktně lineární a~je významně ovlivněn svalovou únavou, která způsobuje pokles střední frekvence (MF) spektra k~nižším hodnotám \cite{Chowdhury2013917, Reaz2006, Stuchla2024}.
Pro kvantitativní hodnocení síly kontrakce se v~časové oblasti standardně využívají příznaky jako integrál absolutních hodnot signálu (iEMG), representující plochu pod křivkou, a~efektivní hodnota (RMS) \cite{Wu2024613, Chowdhury2013917}.
Statistická distribuce amplitudy sEMG přitom vykazuje ne-gaussovské rozdělení při nízkých a~vysokých úrovních síly, zatímco v~středním rozsahu se přibližuje gaussovskému průběhu \cite{Chowdhury2013917}.

Frekvenční spektrum sEMG signálu se pohybuje v~rozsahu 0 až 500~Hz, přičemž dominantní energetická složka se nachází v~pásmu 20 až 150~Hz \cite{Reaz2006, Sul2025627, Wu2024613}.
Složky pod 20~Hz jsou zpravidla zatíženy pohybovými artefakty a~kvazi-náhodnou fluktuací frekvence impulsů motorických jednotek, a~jsou proto při zpracování pro účely řízení eliminovány dolní mezí pásmové propusti \cite{Chowdhury2013917, DeLuca2010, Reaz2006}.

Kvalita zachyceného signálu je kriticky závislá na přesném umístění elektrod nad svalové bříško; signál je zároveň ovlivněn hloubkou svalu a~tloušťkou podkožní tukové vrstvy, která působí jako frekvenčně závislý filtr snižující amplitudu a~zkreslující spektrální složení \cite{ALCAN20231026, Chowdhury2013917, Reaz2006}.
Významnou komplikaci představují přeslechy (\textit{crosstalk}), tedy nechtěné snímání aktivity sousedních svalových skupin, což v~teleoperačních systémech vede k~chybné interpretaci pohybového záměru operátora \cite{Chowdhury2013917, DeLuca2010, Sul2025627}.
Pro kvantifikaci rizika přeslechů se v~literatuře doporučuje sledování metriky \textit{crosstalk risk factor}, jejíž hodnota udává míru pravděpodobnosti kontaminace signálu aktivitou okolních svalů \cite{Wu2024613}.

\subsection{Charakteristiky a omezení sEMG pro řízení}
Povrchový elektromyografický signál je v~odborné literatuře definován jako stochastický proces vznikající časoprostorovou superpozicí akčních potenciálů motorických jednotek (MUAP), jehož amplituda se pohybuje v~řádu mikrovoltů a~frekvenční spektrum je soustředěno v~pásmu 0 až 500~Hz, přičemž dominantní energetická složka se nachází v~rozsahu 20 až 150~Hz \cite{Chowdhury2013917, Reaz2006, Song2023823, Sul2025627}.
Signál vykazuje přibližně gaussovské rozdělení se středem v~nule a~podléhá výrazné inter-individuální variabilitě; svalová únava se projevuje spektrálním posunem k~nižším frekvencím a~nelineárními změnami amplitudy \cite{ALCAN20231026, Lee20241121, Wu2024613}.

Zásadní omezení pro stabilitu řízení představují parazitní složky.
Nízkofrekvenční pohybové artefakty v~pásmu 0–20~Hz jsou způsobeny nestabilitou kontaktu elektrod a~pohyby kabeláže, síťové rušení na frekvenci 50/60~Hz kontaminuje relevantní část spektra a~přeslechy (\textit{crosstalk}), tedy snímání aktivity sousedních svalových skupin, znemožňují přesnou diferenciaci pohybového záměru v~systémech s~nízkým počtem kanálů \cite{ALCAN20231026, Chowdhury2013917, DeLuca2010, Song2023823, Sul2025627}.
K~potlačení těchto vlivů se standardně využívá pásmová propust v~rozsahu 20–500~Hz, jejíž dolní mezní frekvence je kritická pro eliminaci mechanického šumu z~kabeláže a~pohybů kůže \cite{DeLuca2010, Song2023823, Stavdahl2011, Sul2025627}.
Následné vyhlazování pro získání obálky signálu se realizuje metodami jako efektivní hodnota RMS (\textit{Root Mean Square}) odrážející svalové úsilí, nebo průměr absolutních hodnot MAV (\textit{Mean Absolute Value}) pro reprezentaci úrovně aktivace \cite{Sul2025627, Wu2024613}.
Vzhledem k~nelineárnímu vztahu mezi silou kontrakce a~amplitudou signálu je pro opakovatelné řízení nezbytná normalizace vůči maximální volní kontrakci (MVC) podle vztahu
%
\begin{equation}
    sEMG_{\mathrm{nor}} = \frac{sEMG_{\mathrm{task}}}{sEMG_{\mathrm{MVC}}} \times 100\,\%,
\end{equation}
%
kde $sEMG_{\mathrm{nor}}$ je normalizovaná hodnota, $sEMG_{\mathrm{task}}$ aktuální signál a~$sEMG_{\mathrm{MVC}}$ kalibrované maximum uživatele \cite{Wu2024613}.

Z~hlediska stupňů volnosti je jednokanálové sEMG omezeno na binární nebo jednoduché proporcionální řízení; nezávislé ovládání jednotlivých prstů vyžaduje vícekanálové systémy a~pokročilé klasifikátory, které však zvyšují výpočetní náročnost a~senzitivitu vůči variabilitě signálu \cite{Song2023823, Sul2025627}.
Výrazným limitem je rovněž inherentní latence systému.
Přestože sEMG umožňuje predikovat pohyb s~náskokem 50–100~ms před svalovou aktivitou, digitální filtrace a~integrační okna vnášejí do smyčky zpoždění; hodnoty nad 100~ms jsou v~robotice obecně považovány za hranici pro plynulé řízení v~reálném čase \cite{Lee20241121, Song2023823, Sul2025627}, zatímco v~teleoperačních úlohách latence nad 300~ms výrazně narušuje intuitivnost interakce a~nutí operátora přecházet na diskontinuální strategii ovládání \cite{Almeida202098, Bolarinwa2022}.

\section{Řízení v uzavřené smyčce}
Řízení v~uzavřené smyčce (\textit{closed-loop}) představuje v~kontextu haptické teleoperace a~biomechanického biofeedbacku fundamentální mechanismus, při němž jsou fyzikální parametry vzdáleného prostředí kontinuálně snímány, zpracovávány a~formou haptické informace přenášeny zpět k~operátorovi \cite{Kos2018, Suman2025, GIUDICI2025}.
Zatímco systémy s~otevřenou smyčkou vysílají řídicí povely bez ověření jejich skutečného dopadu na okolí, bilaterální teleoperační řetězec uzavírá regulační obvod výměnou energie a~dat mezi stranou operátora a~vzdáleného robotického systému, čímž dosahuje vysoké míry imerze a~přesnosti interakce \cite{Aijaz2018, GIUDICI2025, Thomas2023110}.
Tato zpětná vazba je pro operátora klíčová, neboť mu dovoluje v~reálném čase adaptovat své chování a~intuitivně korigovat úchopovou sílu či trajektorii pohybu na základě vnímaného odporu vzdálených objektů \cite{Mayer2020427, Suman2025, Thomas2023110}.

V~rámci konceptu \textit{Human-in-the-loop} se lidský subjekt stává integrální součástí řídicí smyčky, kde jeho kognitivní funkce a~motorické reakce přímo uzavírají globální cyklus řízení \cite{Aijaz2018, Thomas2023110, Zhang2021530}.
Interaktivní řízení operátora lze formálně popsat vztahem
%
\begin{equation}
    u_{ref} = \frac{1}{\Delta t}\left(x_{2,d} - x_2\right),
\end{equation}
%
kde $u_{ref}$ představuje referenční řídicí vstup operátora, $\Delta t$ je perioda řídicí smyčky, $x_{2,d}$ žádaná a~$x_2$ aktuální rychlost systému \cite{Zhang2021530}.
Výsledná haptická odezva $F$ přenášená zpět k~operátorovi pak může být modelována jako
%
\begin{equation}
    F = K_f \left(u_{CBF} - u_{ref}\right),
\end{equation}
%
kde $K_f$ je konstanta zesílení zpětné vazby a~$u_{CBF}$ představuje bezpečný řídicí signál vypočtený systémem \cite{Zhang2021530}.
Zapojení člověka do regulační smyčky přináší specifické výzvy pro návrh systému, zejména s~ohledem na potřebu zajištění transparence a~stability v~přítomnosti latencí a~jitteru, které nevyhnutelně doprovázejí bezdrátovou komunikaci a~digitální zpracování signálů \cite{Bolarinwa2022, Papakonstantinou20251218, Suman2025}.
Následující podkapitoly se detailně zaměří na teoretické vymezení systémů Human-in-the-loop a~na specifika realizace zpětné vazby síly v~teleoperačních řetězcích.

\subsection{Human-in-the-loop systémy}
Koncept \textit{Human-in-the-loop} (HITL) definuje systémy, v~nichž je lidský operátor integrován jako aktivní prvek v~globální řídicí smyčce, kde dochází k~bidirekcionální výměně energie a~haptických informací mezi řídicí (\textit{master}) a~řízenou (\textit{slave}) doménou \cite{Aijaz2018, Xiong2024}.
Na rozdíl od plně autonomního řízení, které spoléhá na strojovou inteligenci, HITL systémy využívají kognitivní schopnosti a~motorické reakce člověka k~uzavření regulačního obvodu, což umožňuje vysokou míru imerze a~přesnosti v~dálkově ovládaných prostředích \cite{Aijaz2018, Thomas2023110, Zhang2021530}.
Operátor v~tomto řetězci funguje jako kontrolér, jehož výkon je limitován neuromotorickou odezvou a~kognitivní zátěží \cite{Suman2025, Thomas2023110}.
Klíčem k~efektivnímu řízení je kontinuální aktualizace vnitřních dopředných modelů (\textit{internal feedforward models}) operátora na základě haptických vjemů, což vede ke zvýšení neurální efektivity -- tedy vztahu mezi úspěšností úlohy a~vynaloženým mentálním úsilím \cite{Thomas2023110, GIUDICI2025}.

Při implementaci sdíleného řízení (\textit{shared control}) dochází k~interakci mezi záměrem operátora a~asistencí systému, což může vést ke specifickým vzorcům chování: harmonické spolupráci (WH~-- \textit{Working in Harmony}), konfliktu v~cílovém směru (CD~-- \textit{Conflict in target Direction}) či konfliktu v~místě zastavení (CP~-- \textit{Conflict in Parking location}) \cite{AlSaadi2023313}.
Tyto konflikty negativně ovlivňují transparenci a~uživatelskou spokojenost, pokud operátor nerozumí zásahům autonomního kontroléru do své trajektorie \cite{Zhang2021530}.
Pro zvýšení shody záměru se využívá haptická zpětná vazba, kterou lze matematicky modelovat vztahem
%
\begin{equation}
    F = K_f \left(u_{CBF} - u_{ref}\right),
\end{equation}
%
kde $F$ je zpětnovazební síla, $K_f$ konstanta zesílení, $u_{ref}$ referenční vstup operátora a~$u_{CBF}$ představuje bezpečný řídicí signál vypočtený systémem pomocí bariérových funkcí \cite{Zhang2021530}.
Zatímco sEMG vazba umožňuje prediktivní řízení informováním o~svalové aktivitě, kinestetická silová vazba je dominantní modalitou pro snížení kognitivní zátěže při manipulaci s~křehkými objekty \cite{Xiong2024, Thomas2023110, GIUDICI2025, Ninu2014}.
Hlavní výzvu pro stabilitu HITL systémů představuje latence a~jitter v~bezdrátových spojích, které mohou v~uzavřené smyčce způsobit nekontrolované oscilace a~vyžadují robustní bezpečnostní mechanismy pro udržení stability nezávislé na zpoždění \cite{Aijaz2018, Suman2025, GIUDICI2025, Black2025}.

\subsection{Zpětná vazba síly v teleoperaci}
Zpětná vazba síly, v~odborné literatuře definovaná jako kinestetická haptika, představuje v~bilaterální teleoperaci klíčový mechanismus pro přenos interakčních sil ze vzdáleného prostředí k~operátorovi prostřednictvím mechanických aktuátorů na řídicím zařízení \cite{GIUDICI2025, Patel2022, Park2022}.
Tato modalita cílí na hloubkové receptory ve svalech, šlachách a~kloubech, čímž uživateli umožňuje vnímat mechanickou impedanci -- tedy tuhost, hmotnost a~prostorový odpor dálkově ovládaného systému \cite{Yang2025, GIUDICI2025, Li2025}.
Informace o~síle může být snímána buď přímo pomocí tenzometrů či rezistorů citlivých na sílu (FSR) v~místě kontaktu, nebo nepřímo prostřednictvím monitorování proudového odběru servomotorů \cite{Mohammadi2024418, Patel2022}.
Zatímco přímé snímání poskytuje vysokou přesnost a~frekvenční odezvu, nepřímá detekce proudem je konstrukčně robustnější, nevyžaduje senzory v~místě dotyku a~je imunní vůči jejich mechanickému poškození; její citlivost je však limitována vnitřním třením a~mechanickým šumem převodů, což může maskovat jemné změny sil při prvotním kontaktu \cite{Mohammadi2024418, Patel2022}.

Ústředním cílem návrhu je transparence, tedy stav, kdy operátor vnímá vzdálenou úlohu bez zkreslení způsobeného technickým zařízením.
Tento stav lze matematicky popsat hybridní maticí vztahující síly ($f$) a~rychlosti ($v$) na straně Leader ($h$) a~Follower ($n$):
%
\begin{equation}
    \begin{bmatrix} f_h \\ -v_n \end{bmatrix} =
    \begin{bmatrix} h_{11} & h_{12} \\ h_{21} & h_{22} \end{bmatrix}
    \begin{bmatrix} v_o \\ f_n \end{bmatrix}
\end{equation}
%
\cite{Black2025, Almeida202098}.
V~reálných aplikacích je transparence omezena latencí komunikačního řetězce; zpoždění nad 50~ms již významně zkresluje percepci tuhosti a~nad 300~ms může vyvolat kritickou nestabilitu systému \cite{Patel2022, Almeida202098, Bolarinwa2022}.
Mapování mezi naměřenou a~reprodukovanou silou se může řídit exponenciálním modelem ve tvaru
%
\begin{equation}
    F_{\mathrm{user}} = \alpha\left(-e^{\beta F_{\mathrm{robot}}} + \gamma\right),
\end{equation}
%
kde $F_{\mathrm{user}}$ je výsledná síla působící na uživatele, $F_{\mathrm{robot}}$ je síla naměřená senzorem a~konstanty $\alpha$, $\beta$, $\gamma$ definují dynamický rozsah a~citlivost systému \cite{Park2022}.

Implementace silové odezvy prokazatelně zvyšuje úspěšnost manipulace s~křehkými objekty a~snižuje kognitivní zátěž uživatele prostřednictvím zvýšení neurální efektivity -- tedy příznivějšího poměru mezi úspěšností úlohy a~vynaloženým mentálním úsilím \cite{Thomas2023110, Park2022, Ninu2014, Chai2022}.
V~případech, kdy je přímá silová vazba limitována stabilitou či hardwarem, se jako doplněk využívají alternativní modality, jako je vibrotaktilní či elektrotaktilní stimulace, které efektivně informují o~okamžiku kontaktu a~počínajícím prokluzu objektu, čímž uzavírají řídicí smyčku i~při absenci absolutní informace o~síle \cite{Ninu2014, Chai2022, Mayer2020427}.
Pro zajištění bezpečnosti jsou v~reálných systémech nezbytné mechanismy saturace sil a~watchdog timeouty pro automatickou deaktivaci při výpadku bezdrátového spojení \cite{Patel2022, GIUDICI2025}.

\section{Real-time komunikace v mechatronických systémech}
Real-time komunikace v~kontextu mechatronických a~teleoperačních systémů není definována pouze rychlostí přenosu dat, ale především časovou predikovatelností a~determinismem, tedy schopností systému zaručit doručení řídicích a~zpětnovazebních informací v~pevně stanoveném časovém rámci \cite{Aijaz2018, Kos2018, Bolarinwa2022}.
V~uzavřených řídicích smyčkách se silovou zpětnou vazbou představuje latence kritický parametr, který přímo ovlivňuje stabilitu systému a~jeho transparenci, tedy míru imerze, se kterou operátor vnímá vzdálené prostředí \cite{Suman2025, Aijaz2018, Bolarinwa2022, Patel2022, Black2025}.
Celkové zpoždění v~komunikačním kanálu $t_{delay}$ lze analyticky popsat vztahem
%
\begin{equation}
    t_{delay} = t_{prop} + t_{tran} + t_{MAC},
\end{equation}
%
kde $t_{prop}$ je doba šíření signálu, $t_{tran}$ představuje čas vlastní transmise rámce a~$t_{MAC}$ odpovídá zpoždění protokolu pro přístup k~přenosovému médiu \cite{Kos2018}.
Čas přenosu je přitom přímo úměrný délce datového rámce $L$ a~nepřímo úměrný bitové rychlosti kanálu $R$ dle vztahu
%
\begin{equation}
    t_{tran} = \frac{L}{R}
\end{equation}
%
\cite{Kos2018}.

Zatímco drátová spojení vykazují vysoký stupeň determinismu, bezdrátový přenos informací představuje technologický kompromis mezi flexibilitou nasazení a~přísnými nároky na reálný čas; technologie jako Bluetooth~Classic nabízejí výhody v~podobě nízké energetické náročnosti, avšak jsou zatíženy inherentní latencí a~rizikem kolísání zpoždění, označovaného jako jitter \cite{Aijaz2018, Eridani2021918, Bolarinwa2022, Suman2025}.
Překročení mezních hodnot latence, které jsou pro plynulou haptickou teleoperaci vnímány jako kritické při dosažení hranice 300 až 400~ms, vede k~narušení stability regulační smyčky a~nezřídka vynucuje změnu strategie řízení na diskontinuální pohyby \cite{Bolarinwa2022, Suman2025, Aijaz2018, AlSaadi2023313, Almeida202098}.
Následující podkapitoly se detailně zaměří na teoretické vymezení latence a~determinismu v~digitálních systémech a~na specifické limity bezdrátové komunikace v~reálných mechatronických aplikacích.
\subsection{Latence a determinismus}
Real-time komunikace v~mechatronických a~teleoperačních systémech je definována především časovou predikovatelností a~determinismem, což představuje schopnost technického řešení zaručit doručení řídicích a~zpětnovazebních informací v~pevně stanoveném časovém rámci \cite{Aijaz2018, Kos2018, Bolarinwa2022}.
V~uzavřených řídicích smyčkách se silovou zpětnou vazbou představuje celková latence kritický parametr, který přímo ovlivňuje stabilitu systému a~jeho transparenci \cite{Patel2022, Suman2025, Black2025}.
Zpoždění v~komunikačním řetězci $t_{delay}$ lze analyticky vyjádřit jako součet doby šíření signálu $t_{prop}$, doby vlastní transmise rámce $t_{tran}$ a~zpoždění protokolu pro přístup k~přenosovému médiu $t_{MAC}$ \cite{Kos2018}.
Přenosové zpoždění je přitom přímo úměrné délce datového paketu $L$ a~nepřímo úměrné bitové rychlosti kanálu $R$ dle vztahu $t_{tran} = L/R$, což činí volbu baudové rychlosti UART a~optimalizaci velikosti paketu klíčovým faktorem pro real-time odezvu \cite{Kos2018}.

V~haptických aplikacích jsou nároky na latenci extrémně přísné; pro věrnou silovou odezvu je kritickou hranicí zpoždění 50~ms, nad nímž dochází k~výraznému zkreslení percepce tuhosti prostředí \cite{Patel2022, Suman2025, Almeida202098}.
Pro oblast robotického a~protetického řízení v~uzavřené smyčce je jako vhodný rozsah celkového zpoždění uváděno 100 až 250~ms, přičemž zpoždění přesahující 300~ms vede k~prokazatelnému poklesu úspěšnosti úloh \cite{Song2023823, Almeida202098}.
Nedeterministické kolísání latence, označované jako jitter, vnáší do bilaterálního řetězce fázový posun, který může vyvolat nestabilitu projevující se oscilacemi nebo vnímáním trhavého pohybu \cite{Patel2022, Black2025}.
Výzkumy přitom naznačují, že u~sledovacích úloh je průměrná hodnota latence pro celkovou výkonnost řízení kritičtější než samotná nahodilost jejího kolísání \cite{Suman2025}.
Překročení hranice 300 až 400~ms pak obvykle vynucuje změnu strategie operátora na diskontinuální režim „pohyb a~čekání" \cite{Suman2025, Bolarinwa2022, AlSaadi2023313, Almeida202098}.

Zdroje těchto zpoždění zahrnují vedle samotné bezdrátové transmise a~Bluetooth stacku také digitalizaci signálů, výpočetně náročnou filtraci dat a~buffering na straně přijímače, který sice kompenzuje jitter, avšak za cenu navýšení latence v~nejhorším případě (\textit{worst-case latency}) \cite{Patel2022, Kos2018, Suman2025}.
Pro dosažení vysoké transparence a~stability je proto v~real-time mechatronice nezbytné minimalizovat $t_{tran}$ zvýšením bitové rychlosti, optimalizovat délku paketu a~využívat technologie s~nízkou inherentní latencí, jako je ESP-NOW nebo optimalizované Bluetooth protokoly \cite{Kos2018, Eridani2021918, Papakonstantinou20251218}.

\subsection{Bezdrátová komunikace a její limity}
Bezdrátové komunikační protokoly v~mechatronice a~teleoperaci představují technologický kompromis mezi dosahem, energetickou náročností a~latencí, přičemž mezi nejrozšířenější standardy patří Bluetooth~Classic, Bluetooth Low Energy (BLE), Wi-Fi a~proprietární protokol ESP-NOW \cite{Eridani2021918, Urazayev202354, Bolarinwa2022}.
Srovnávací studie ukazují, že Bluetooth~Classic vyniká především nízkou spotřebou energie (přibližně 141~mW v~režimu přijímače), avšak zaostává v~přenosové latenci, která se u~1B paketu pohybuje kolem 6~ms, a~v~maximálním dosahu v~podmínkách přímé viditelnosti, jenž činí zhruba 15~metrů \cite{Eridani2021918}.
Naproti tomu protokol ESP-NOW, běžící nad fyzickou vrstvou IEEE~802.11, dosahuje v~identických podmínkách latence pouze 1~ms a~dosahu až 185~metrů s~integrovanou anténou, což z~něj činí vysoce responzivní řešení pro řídicí systémy \cite{Eridani2021918, Urazayev202354, Papakonstantinou20251218}.
Wi-Fi (standard 802.11n) nabízí nejvyšší datovou propustnost v~řádu desítek~Mbps, avšak jeho energetická náročnost a~latence (kolem 3,3~ms) jsou ve srovnání s~ESP-NOW vyšší \cite{Eridani2021918, Urazayev202354}.

Specifické limity bezdrátového přenosu pro haptické aplikace zahrnují zejména proměnné zpoždění (jitter), výpadky paketů a~vliv fyzických překážek, které mohou u~frekvence 2,4~GHz dramaticky snížit kvalitu signálu již po průchodu dvěma až třemi stěnami \cite{Suman2025, Urazayev202354}.
Výzkumy naznačují, že zatímco lidský nervosvalový systém dokáže do jisté míry tolerovat náhodné fluktuace latence, průměrná hodnota zpoždění je kritickým faktorem pro výkonnost řízení; výrazný pokles kvality teleoperace nastává při překročení hranice 400~ms \cite{Suman2025, Bolarinwa2022}.
Budoucí standardy 5G a~6G v~rámci ultra-spolehlivé komunikace s~nízkou latencí (URLLC) cílí na dosažení doby zpátečního přenosu (\textit{round-trip time}, RTT) na úrovni 1~ms a~spolehlivost doručení $10^{-5}$, což je nezbytné pro plnou imerzi v~rámci taktilního internetu \cite{Awais202375, Aijaz2018}.

Pro zajištění spolehlivosti v~reálných podmínkách jsou nezbytné softwarové strategie, jako jsou watchdog mechanismy a~definované \textit{fail-safe} stavy, které při ztrátě spojení uvedou systém do bezpečného režimu \cite{Bolarinwa2022, Black2025}.
Nestabilita v~haptické smyčce sice u~lidského operátora nepředstavuje přímé bezpečnostní riziko jako u~průmyslových robotů, avšak vede k~znemožnění dálkového vedení pohybu \cite{Black2025}.
Modelování ztrát paketů se často realizuje pomocí Gilbert-Elliottova modelu, kde průměrná chybovost $p_E$ je dána poměrem pravděpodobností přechodu mezi dobrým a~špatným stavem kanálu jako
%
\begin{equation}
    p_E = \frac{u}{u + v}
\end{equation}
%
\cite{Suman2025}.
Implementace robustních komunikačních vrstev je proto klíčem k~zachování transparence a~stability bilaterální teleoperace v~nedeterministickém bezdrátovém prostředí \cite{Aijaz2018, Bolarinwa2022, Patel2022, Kos2018}.

\section{Kinematika ruky a mapování pohybu}
Analýza kinematiky lidské ruky představuje klíčový teoretický základ pro návrh a~realizaci moderních teleoperačních systémů, neboť umožňuje definovat přesné geometrické a~mechanické vztahy potřebné pro věrný přenos pohybového záměru operátora na robotický efektor \cite{Belcamino2024, GIUDICI2025, Mizrahi2023}.
Ústřední výzvu v~této oblasti tvoří výrazná biomechanická diskrepance mezi lidskou rukou, která je v~závislosti na zvolené komplexitě modelu popisována 21 až 27~stupni volnosti (DOF), a~robotickými systémy, jež jsou z~konstrukčních důvodů často redukovány na 5 až 6~nezávisle řízených kanálů \cite{Belcamino2024, Devaraja20201222, Zhang2025}.
Tento biomechanický nesoulad, označovaný jako \textit{embodiment gap}, vyžaduje nasazení algoritmů pro retargeting pohybu \cite{Zhang2025}.
Retargeting je proces transformace zachycených pohybových dat z~ovládacího zařízení do konfiguračního prostoru robotického efektoru tak, aby byla zachována intuitivnost a~funkčnost úchopů i~při omezené kinematické věrnosti technického zařízení \cite{GIUDICI2025, Zhang2025}.

V~základní formě lze toto mapování modelovat pomocí lineární transformace
%
\begin{equation}
    \theta_{follower} = s \cdot \theta_{leader} + o,
\end{equation}
%
kde $\theta_{follower}$ je cílový úhel servomotoru robotické ruky, $\theta_{leader}$ je vstupní hodnota ze snímače polohy operátora, parametr $s$ definuje měřítko pohybu (\textit{scaling factor}) a~$o$~reprezentuje kalibrovaný posun (\textit{offset}) pro dosažení shody v~krajních polohách \cite{Zhang2025, Mizrahi2023}.
Následující podkapitoly se detailně zaměří na systematické srovnání stupňů volnosti u~biologických a~technických systémů a~na konkrétní principy retargetingu, které umožňují plynulé a~efektivní mapování pohybu v~rámci bilaterálních teleoperačních řetězců.

\subsection{Stupně volnosti lidské a robotické ruky}
Lidská ruka představuje z~biomechanického hlediska vysoce komplexní mechanismus, který v~závislosti na zvolené úrovni abstrakce a~zahrnutí zápěstí disponuje 21 až 27~stupni volnosti (DOF) \cite{Belcamino2024, Devaraja20201222, GIUDICI2025}.
Samotné prsty zahrnují 21 volně ovladatelných stupňů volnosti, kdy každý ze čtyř prstů disponuje čtyřmi DOF (dva v~metakarpofalangálním kloubu pro flexi-extenzi a~addukci-abdukci, jeden v~proximálním a~jeden v~distálním interfalangálním kloubu) a~palec využívá pět DOF díky specifickému trapeziometakarpálnímu kloubu \cite{Belcamino2024, Capek2018, Cihak2011}.
Zbývajících 6 DOF je lokalizováno v~zápěstí a~umožňuje flexi-extenzi, addukci-abdukci, pronaci-supinaci a~translaci v~prostoru \cite{Belcamino2024}.
Oproti této biologické redundanci jsou robotické systémy z~konstrukčních důvodů často redukovány na mnohem nižší počet nezávisle řízených kanálů, typicky 5 až 16 -- například model InMoov s~pěti servo-poháněnými prsty -- což vytváří takzvaný \textit{embodiment gap}, tedy biomechanický nesoulad mezi operátorem a~technickým zařízením \cite{Devaraja20201222, GIUDICI2025, Zhang2025}.

Lidský mozek zvládá vysokou dimenzionalitu pohybu prostřednictvím motorických synergií, které umožňují koordinovat pohyb prstů jako funkčních celků místo individuálního řízení každého kloubu \cite{Mizrahi2023, Xiong2024, Zhang2025}.
Omezení DOF u~robotických rukou pro operátora znamená kritickou ztrátu kinematické redundance, což komplikuje manipulaci v~kontaktu s~objekty a~omezuje schopnost jemného polohování \cite{Belcamino2024, GIUDICI2025}.
Věrnost retargetingu pohybu je přitom přímo závislá na počtu DOF sledovacího systému; zatímco jednoduchá senzorická rozhraní snímající pouze 5~DOF jsou vhodná pro základní gesta a~binární úchopy, pokročilé systémy schopné zachytit plný rozsah lidské kinematiky výrazně zvyšují plynulost a~přesnost dálkového ovládání \cite{Belcamino2024, Zhang2025}.
Pro kompenzaci chybějících DOF se využívá lineární transformace
%
\begin{equation}
    \theta_{follower} = s \cdot \theta_{leader} + o,
\end{equation}
%
kde $\theta_{follower}$ je vypočtený úhel pro servomotor robota, $\theta_{leader}$ vstupní hodnota ze senzoru operátora a~parametry $s$ (měřítko) a~$o$ (posun) definují funkční transformaci mezi konfiguračními prostory obou zařízení \cite{Zhang2025, Mizrahi2023, GIUDICI2025}.
Přesné snímání, například pomocí sEMG či exoskeletů, v~kombinaci s~těmito algoritmy umožňuje operátorovi udržet intuitivní kontrolu i~nad kinematicky chudším systémem \cite{Xiong2024, Mizrahi2023, GIUDICI2025}.

\subsection{Principy retargetingu pohybu}
Motion retargeting v~kontextu haptických teleoperačních systémů představuje proces transformace zachycených pohybových dat z~konfiguračního prostoru lidské ruky (zdrojová doména) do prostoru robotického efektoru (cílová doména) tak, aby byla zachována funkčnost a~intuitivnost ovládání i~při výrazných kinematických rozdílech \cite{GIUDICI2025, Zhang2025}.
Tento proces je nezbytný pro překonání \textit{embodiment gap}, tedy biomechanického nesouladu vyplývajícího z~odlišné geometrie kloubů, délky kostních segmentů a~rozsahu pohybu (ROM) obou systémů \cite{Coppola2022920, GIUDICI2025, Zhang2025}.
V~základní implementaci se využívá lineární mapování polohy vyjádřené vztahem
%
\begin{equation}
    \theta_{follower} = s \cdot \theta_{leader} + o,
\end{equation}
%
kde $\theta_{follower}$ je cílový úhel servomotoru robota, $\theta_{leader}$ vstupní hodnota ze senzoru (např.\ potenciometru), $s$ definuje měřítko pohybu (\textit{scaling factor}) a~$o$~představuje kalibrovaný posun (\textit{offset}) pro sjednocení krajních poloh \cite{Mizrahi2023, Zhang2025}.
Pokročilejší metody pracující v~kartézském prostoru využívají vážený průměr polohového a~rychlostního mapování
%
\begin{equation}
    \hat{R}_f = \beta \cdot \hat{R}_{f,pos} + (1-\beta) \cdot \hat{R}_{f,vel},
\end{equation}
%
kde $\hat{R}_f$ představuje výslednou pozici konečku prstu, koeficient $\beta$ váží příspěvky obou složek a~škálovací koeficienty dále kompenzují strukturální rozdíly v~délce článků prstů a~dlaní obou systémů \cite{Coppola2022920}.

Při výrazné neshodě v~počtu stupňů volnosti -- kdy lidská ruka disponuje až 27~DOF oproti typickým 5--6~DOF u~levnějších robotických náhrad -- se retargeting opírá o~motorické synergie \cite{Belcamino2024, GIUDICI2025, Mizrahi2023, Xiong2024}.
Ty umožňují redukovat komplexitu řízení shlukováním pohybů do funkčních celků, přičemž je nutné zohledňovat i~případné rozdíly v~pořadí a~orientaci prstů v~komunikační vrstvě \cite{Mizrahi2023, Zhang2025}.
Specifickou výzvu představuje retargeting v~myoelektrických systémech; zatímco vícekanálové sEMG systémy využívají hluboké učení -- například architektury MSCNN (\textit{Multi-stream Convolutional Neural Network}) nebo LSTM (\textit{Long Short-Term Memory}) -- pro mapování svalové aktivity na plynulý pohyb prstů, jednokanálové systémy jsou inherentně omezeny na binární či jednoduché proporcionální řízení, kde chybí nezávislá informace o~aktivaci konkrétních svalových skupin pro každý prst zvlášť \cite{Mizrahi2023, Sun2023, Xiong2024}.
Zatímco jednokanálové sEMG řízení tak vytváří přirozené omezení retargetingu, věrné polohové snímání prokazatelně zvyšuje úspěšnost úloh, minimalizuje prokluz objektů a~snižuje kognitivní zátěž operátora prostřednictvím zvýšení neurální efektivity, přičemž kvalita retargetingu -- zahrnující správnou orientaci a~synchronizaci prstů -- přímo determinuje míru imerze ve vzdáleném prostředí \cite{Coppola2022920, Li2025, Thomas2023110, Zhang2025}.
